\section{Higher-moment and directional consistency arguments}
\label{app:directional}
For Proposition~\ref{prop:above-two}, conditional versions of
\eqref{eq:utility-moment-conversion} hold whenever the stated return-moment
contract holds conditionally. Write $A_p$ for that public bound and
$\kappa_p=(p-1)^{p-1}/p^p$. Choose $B_i=B_0 2^i$, $i\ge0$, $B_0>0$, and
\[
 b_i=\kappa_p A_p/B_i^{p-1},\qquad
 X_{i,t}=\max\{U(r_t),-B_i\}-b_i,\qquad
 \lambda_{ij}=\frac{2^{-j-1}}{B_i+b_i},\quad i,j\ge0.
\]
The sharp tail inequality \eqref{eq:sharp-tail} gives
$\E_{t-1}X_{i,t}\le\E_{t-1}U(r_t)$. Each factor
$1+\lambda_{ij}X_{i,t}\ge1-2^{-j-1}\ge1/2$ therefore has conditional
expectation at most one under the original utility null. With weights
$\omega_{ij}=2^{-i-j-2}$, the sum of the products is a valid e-process by
conditional Tonelli. Its finite-time value is finite on every finite return
path: $\lambda_{ij}\le1/(2B_0)$ and each factor is at most
$1+|U(r_t)|/(2B_0)$, giving a common finite product upper bound.

Under an iid law with $\delta=\E U(r)>0$, select a fixed $i$ with
$b_i<\delta/2$. Then $\E X_i\ge\delta-b_i>0$. For a fixed integrable $X$
bounded below by $-D$, dominated convergence gives
\begin{equation}
 \frac{\E\log(1+\lambda X)}{\lambda}\longrightarrow\E X
 \quad\text{as }\lambda\downarrow0,\quad \lambda D\le1/2.
 \label{eq:small-stake-growth}
\end{equation}
For $X\ge0$ the quotient is at most $X$; for $-D\le X<0$ its absolute value
is at most $2|X|$. These bounds prove domination and log integrability.
Some $j$ thus has positive expected log factor. The iid strong law makes
that component grow to infinity almost surely. Its fixed positive mixture
weight proves eventual crossing at every individual level, completing
Proposition~\ref{prop:above-two}. No finite expectation of the hitting time
is needed or inferred.

For the cash-PARK mixture, $Y_t\ge-1/(2\gamma)$ makes every dyadic linear
factor positive and of null conditional expectation at most one. The weights
sum to one, so the same conditional-Tonelli argument proves validity without
a supplied numerical return-moment bound. With iid integrable positive-mean
$Y$, \eqref{eq:small-stake-growth} supplies a fixed component with positive
log growth and proves eventual power one. This assertion does not apply
the iid strong law to the dependent estimated-policy scores of
Theorem~\ref{thm:vanishing}; that theorem has its own proof below.
